\Hefteintrag{1.5}{Besondere Methoden: Konstruktoren}{%
Konstruktoren sind spezielle Methoden, die genutzt werden, um ein \emphColA{Objekt zu initialisieren/konstruieren}. Ein Konstruktor wird \emphColA{automatisch beim Erstellen eines Objekts} aufgerufen. 

In Java erkennt man einen Konstruktor insb. an zwei Dingen:
\begin{enumerate}
    \item Er \emphColA{heißt} \emphColA{\LoesungLuecke{genauso wie die Klasse}{12cm}}.
    \item Er hat \emphColA{keinen \LoesungLuecke{Rückgabetyp}{8cm}}.
\end{enumerate}

Ansonsten verhält er sich wie eine Methode und kann z.B. Parameter haben, die dann oft als Startwerte für das Objekt dienen.

Ein typischer Konstruktor im Programmcode sieht z.B. so aus: \\
\emphColB{public Spieler(double startX , double startY) \{ \\ 
~~~~~~x = startX;\\
~~~~~~y = startY;\\
\}}

Im Klassendiagramm würde dieses Beispiel so eingetragen:

\vspace{0.2cm}

\emphColB{\LoesungLuecke{+ Spieler(double startX, double startY)}{12cm}}




Ein neues Objekt erstellt und speichert man dann z.B. so: \emph{Spieler s1 = new Spieler(100,50);}




}